\documentclass[12pt,a4paper]{article}
\usepackage[margin=1in]{geometry}
\usepackage{amsmath,amssymb}
\usepackage{hyperref}

\title{\textbf{The Cascade Series}}
\author{RTAC}
\date{March 2026}

\begin{document}
\maketitle
\thispagestyle{empty}

\section*{The Thought Experiment}

A black hole has a horizon. An infalling observer sees infinite time dilation there.
But the black hole evaporates in finite time. These two facts are in tension.

The resolution is asymptotic: there is no sharp horizon. Infalling matter never
crosses --- from outside, it asymptotically approaches the horizon; from the infaller's
perspective, the black hole evaporates and the horizon recedes. Both observers agree:
the infaller asymptotically tracks a shrinking shell where the rate of infall balances
the rate of evaporation.

This shell is the physical content of the ``horizon.'' On it sits a two-dimensional
spatial surface ($S^2$ for a Schwarzschild hole), and the time direction is the asymptotic
balance itself --- the direction along which infall and shrinkage compete. What if
this shell is a $2 + 1$ dimensional universe?

If a 4D black hole produces a $2 + 1$ shell universe, does a 5D black hole produce
a $3 + 1$ shell universe? Is that us? Are we living on the $S^3$ horizon of a 5D black
hole? If so, time is the forced direction --- the compactification that the cascade
must undergo, not a coordinate we choose. Moving orthogonally, along the shell,
slows progress in this forced direction but can never stop it completely. This is time
dilation, from the geometry of a shell.

If yes, then the 5D black hole sits inside a 6D space --- another shell of a 7D object.
Follow the tower upward: each shell is the horizon of the next. The tower terminates
at $d = \infty$: the infinite-dimensional unit ball, which has zero volume, zero surface
area, no interior. Geometric nothing. A natural starting point.

That was the end of the intuition. The question was: what happens if you work
back down from infinity?

\section*{The Starting Point}

The thought experiment arrives at the infinite-dimensional unit ball from above,
by following a tower of black hole shells upward. But there is a deeper reason
to start there, established in the Prelude: \emph{there is nowhere else to start.}

A theory of everything cannot have inputs. Every input demands an explanation of
its origin, and that explanation either requires a deeper theory or is circular.
The only starting point that requires no explanation is nothing.

Nothing --- taken seriously as a mathematical object --- is not featureless.
The precondition for any formal system is the distinction between true and false:
$0\neq 1$. Distinction implies orthogonality (zero mutual information between
distinguishable states). Orthogonality iterates without bound (no logical
obstruction to countable independence). The absence of external scale forces unit
norm. The result is the infinite-dimensional unit ball: zero volume, zero
surface area, no interior. Structured nothing.

The unit ball is not chosen from a menu of starting points. It is the unique
mathematical structure that corresponds to ``nothing with the capacity for
distinction,'' derived from the precondition for coherent thought. The series
is the derivation of what that structure implies.

\section*{The Hypothesis}

\begin{center}
\textbf{The infinite-dimensional unit ball, descended to four dimensions, is indistinguishable from our universe.}
\end{center}

It is tested by deriving, from the cascade's geometry alone, the cosmological
constant, the dimension and signature of spacetime, quantum mechanics, and the
Standard Model gauge group and its symmetry breaking, three fermion generations,
fermion masses, and gauge coupling constants. It further derives the background
cosmological parameters --- including a Hubble constant between the two competing
measurements and an account of the DESI baryon acoustic oscillation observations
without dynamical dark energy. Every prediction is a test of the hypothesis.

\section*{The Series}

\textbf{Prelude: Why Nothing Has Structure.} From $0\neq 1$ to $B^\infty$
in five steps: distinction $\to$ orthogonality $\to$ infinite dimensions $\to$
unit norm $\to$ the infinite-dimensional unit ball. No step introduces a free
parameter. No step selects from alternatives. The starting point is established.

\textbf{Part 0: Scale Variance from Orthogonality.} Pure mathematics. The slicing
recurrence of the unit ball contains one constant ($\sqrt{\pi}$), whose natural zero
generates two thresholds at $d = 19$ and $d = 217$. The Gamma function produces exactly
four distinguished dimensions: the volume maximum ($d = 5$), the area maximum ($d = 7$),
and the two thresholds. No fifth exists. The cascade invariant
$\Omega_{19}\times\Omega_{217} = 1.205\times 10^{-120}$ is forced. Every step is a theorem
about the Gamma function. No physics enters.

\textbf{Part I: The Cosmological Constant from the Observer's Frame.} The hypothesis
enters. We observe four dimensions. The observer at $d = 4$ lives on the $S^3$ boundary
of $d = 5$, the volume maximum. The frame correction from the observer's host to the
cascade's natural reference is $(\Omega_5/\Omega_7)^2 = 9/\pi^2$:
\[
\frac{\rho_\Lambda}{M^4_{\mathrm{Pl}}}
= \frac{9\;\Omega_{19}\;\Omega_{217}}{\pi^2}
= 1.099\times 10^{-120}.
\]
Observed: $(1.10\pm 0.02)\times 10^{-120}$. Deviation: $0.1\%$. The vacuum is 198
discrete layers of geometry with a natural floor at $\Omega_{217}\approx 10^{-120}$.

\textbf{Part II: Quantum Mechanics from the Cascade.} Consecutive slicing axes are forced
orthogonal, producing complex structure, Hilbert space, the Born rule, and unitary
evolution --- without quantum postulates.

\textbf{Part III: General Relativity from the Cascade.} The cascade's complex spinor
structure intersected with Lovelock uniqueness forces $d = 4$, Lorentzian signature, and
Einstein's equation. The dark energy equation of state is $w = -1$ exactly.

\textbf{Part II\,=\,III: Quantum Gravity without Quantising Gravity.} The quantum and
gravitational projections share the same source, propagator, state space, and hypothesis.
Gleason forces the Born rule; Lovelock forces Einstein's equation; both are unique at
$d = 4$. The standard QM/GR conflicts dissolve. Black hole entropy $S = A/4$ is hidden
geometry from boundary dominance: the factor $1/4$ is $1/d$ at the observer's dimension.

\textbf{Part IVa: The Standard Model from the Cascade (Gauge Group).} Bott periodicity,
Adams' theorem, and the hairy ball theorem jointly force
$\mathrm{SU}(3)\times\mathrm{SU}(2)\times\mathrm{U}(1)$, its symmetry breaking pattern,
and three fermion generations.

\textbf{Part IVb: The Standard Model from the Cascade (Masses and Couplings).} The
geometric-topological factorisation of the fermion mass gives fourteen precision
predictions --- all sub-3\%, zero free parameters.

\textbf{Part V: Cosmology from the Cascade.} The cascade derives $\Omega_m = 1/\pi$,
$\Omega_b = 1/(2\pi^2)$, $H_0 = 71.05$~km/s/Mpc, and $T_{\mathrm{CMB}} = 2.730$~K. The
Friedmann equation has every coefficient determined by $\pi$. The cascade's
$r_d\approx 141$~Mpc fits DESI DR2 BAO with $w = -1$ throughout; the apparent DESI
preference for $w\neq -1$ is a zero-parameter consequence of imposing the wrong ruler.

\bigskip

The hypothesis is never falsified. Papers I--III and II\,=\,III are exact: every result
matches observation with no approximation and no discrepancy. Papers IV--V use
leading-order approximations; their predictions separate into two populations:
descent-dependent quantities carry uniformly negative deviations of $-0.13\%$ to
$-2.20\%$; geometric quantities carry positive deviations of $+0.16\%$ to $+2.76\%$.
The clean sign separation identifies the subleading cascade descent as the source of
the dominant systematic, not computed in this series. Falsification requires completing
an approximate derivation to full exactness and finding a result that does not fit.

\bigskip

The thought experiment that opens this page is the physical content of the series'
central result. The observer is on the $S^3$ shell of a 5D black hole, partway through
an asymptotic compactification that never completes. The cosmological constant is the
cascade's geometry measured from that shell. Time dilation is the local rate of the
compactification. Black hole horizons are where the transition locally completes. The
cosmological constant and time dilation are proportional, in the ratio $d = 5$. The
Gamma function knows how big the universe is because the universe is the unit ball,
descended.

\end{document}